\begin{sgastatement}{\SGAInlineRefTarget{sga3:xxvi:corollary:5-2}Corollary 5.2}
Let \(S\) be a semilocal scheme, \(G\) an \(S\)-reductive group, and
\(P,P'\) two opposite parabolic subgroups of \(G\).  The canonical map
\[
 \rad^u(P)(S)\cdot\rad^u(P')(S)\longrightarrow(G/P)(S)
\]
is surjective.  In particular,
\[
 (G/P)(S)=G(S)/P(S).
\]
Every parabolic subgroup \(Q\) of \(G\) having the same type as \(P\)
is of the form
\[
 Q=\Int(uu')P
\]
with \(u\in\rad^u(P)(S)\) and \(u'\in\rad^u(P')(S)\).
\end{sgastatement}
\SGARefTarget{sga3:xxvi:cor:5:2}

The second assertion is plainly equivalent to the first, so let us
prove the latter.  Let \(s_i\) be the closed points of \(S\), let
\(V\) be the open subscheme of \(G/P\) which is the image of \(P'\)
\textup{(and is isomorphic to \(\rad^u(P')\), by \SGARefLink{sga3:xxvi:remark:4:3:6}{4.3.6})}, and take
\(x\in(G/P)(S)\).  By \SGARefLink{sga3:xxvi:theorem:5:1}{5.1}, for every \(i\) there is
\[
 u_i\in\rad^u(P)\bigl(\kappa(s_i)\bigr)
\]
such that \(u_ix_{s_i}\) is a section of \(V_{s_i}\).  Choose
\(u\in\rad^u(P)(S)\) lifting all the \(u_i\), as permitted by \SGARefLink{sga3:xxvi:cor:2:6}{2.6}.
Then \(ux\) is a section of \(V\), since this may be checked on the
closed fibers.  Finally,
\[
 \rad^u(P')(S)\xrightarrow{\sim}V(S)
\]
is bijective, and the result follows.

\begin{sgastatement}{Corollary 5.3}
Let \(S\) be a semilocal scheme, \(G\) an \(S\)-reductive group, and
\(P,P',Q\) three parabolic subgroups of \(G\).  There is a
\(g\in G(S)\) such that \(\Int(g)Q\) is in transverse position
relative to both \(P\) and \(P'\).
\end{sgastatement}
\SGARefTarget{sga3:xxvi:cor:5:3}

In the notation of \SGARefLink{sga3:xxvi:corollary:4-2-4}{4.2.4}\textup{(ii)}, we must show that the
universally schematically dense open subscheme
\[
 \Gen(Q/P)\cap\Gen(Q/P')
\]
of \(G\) has a section over \(S\).
Choose a parabolic subgroup \(Q_1\) opposite to \(Q\), by
\SGARefLink{sga3:xxvi:corollary:4-3-5}{4.3.5}\textup{(i)}, and put
\[
 U=\rad^u(Q),\qquad U'=\rad^u(Q_1).
\]
We shall find the desired \(g\) in \(U(S)U'(S)\).  In this form,
\SGARefLink{sga3:xxvi:corollary:4-2-4}{4.2.4}\textup{(i)} and \SGARefLink{sga3:xxvi:cor:2:6}{2.6} show that it suffices to work on the fibers
over the closed points of \(S\); we may therefore suppose
\(S=\operatorname{Spec}(k)\).

If \(k\) is algebraically closed, there is a \(g\in G(k)\) with the
required property.  By \SGARefLink{sga3:xxvi:corollary:5-2}{5.2} it may be written
\[
 g=uu'q,\qquad u\in U(k),\quad u'\in U'(k),\quad q\in Q(k),
\]
and \(\Int(uu')Q=\Int(g)Q\).

Suppose \(k\) is infinite.  Let \(V\) be the open subset of
\(U\times_kU'\) defined by the cartesian diagram
\[
\begin{tikzcd}[column sep=huge,row sep=large]
 U\times_kU' \arrow[r] & G\\
 V \arrow[u,hook] \arrow[r]
   & \Gen(Q/P)\cap\Gen(Q/P') \arrow[u,hook].
\end{tikzcd}
\]
Since \(V(\bar k)\ne\varnothing\), as just shown, \(V\) is dense in
\(U\times_kU'\); it therefore has a \(k\)-point by \SGARefLink{sga3:xxvi:cor:2:7}{2.7}.

Suppose finally that \(k\) is finite.  Lang's theorem, as used in
5.1.5, shows that \(P\), respectively \(P'\), has a Borel subgroup
\(B\), respectively \(B'\), since
\[
 \Bor(P)\simeq\Bor(P/\rad^u(P)),\qquad
 \Bor(P')\simeq\Bor(P'/\rad^u(P'))
\]
are smooth.
\footnote{\emph{Translator's note.} The displayed French source prints
\(\Bor(Q)\) in the second isomorphism; the parallel argument requires
\(\Bor(P')\), which is used here.}
Let \(B_1\) be a Borel subgroup opposite to \(B\).  By \SGARefLink{sga3:xxvi:corollary:5-2}{5.2} there are
\[
 a\in B^u(k),\qquad a_1\in B_1^u(k)
\]
such that \(\Int(aa_1)B=B'\).  Then
\[
 B_0=\Int(aa_1)B_1=\Int(a)B_1
\]
is opposite to both \(B'\) and \(B\).  Let \(Q_0\) be the unique
parabolic subgroup of \(G\) containing \(B_0\) and having the same
type as \(Q\), as in \SGARefLink{sga3:xxvi:lemma:3:8}{3.8}.  It is in transverse position relative to
\(P\) and \(P'\), by \SGARefLink{sga3:xxvi:theorem:4-2-1}{4.2.1}\textup{\SGARefLink{sga3:xxvi:theorem:4-2-1:item:vi}{(vi)}}.  On the other hand, \SGARefLink{sga3:xxvi:corollary:5-2}{5.2}
gives
\[
 Q_0=\Int(uu')Q,\qquad u\in U(k),\quad u'\in U'(k),
\]
which proves the assertion.

\begin{sgastatement}{Corollary 5.4}
Let \(S\) be a semilocal scheme, \(G\) an \(S\)-reductive group, and
\(P,Q\) two parabolic subgroups of \(G\).  There is a \(g\in G(S)\)
such that \(\Int(g)P\) is in osculating position relative to \(Q\);
equivalently, \(\Int(g)P\cap Q\) is a parabolic subgroup of \(G\).
\end{sgastatement}
\SGARefTarget{sga3:xxvi:cor:5:4}

Choose a parabolic subgroup \(P'\) opposite to \(P\), by
\SGARefLink{sga3:xxvi:corollary:4-3-5}{4.3.5}\textup{(i)}.  By \SGARefLink{sga3:xxvi:cor:5:3}{5.3}, there is a parabolic subgroup \(P'_1\) of
the same type as \(P'\), in transverse position relative to both
\(P\) and \(Q\).  Let \(T\) be a maximal torus of \(P'_1\cap Q\),
as in \SGARefLink{sga3:xxvi:proposition:4-2-7}{4.2.7}\textup{(ii)}, and let \(P_1\) be the parabolic subgroup
opposite to \(P'_1\) relative to \(T\).  Then \(P_1\) and \(Q\) are
in osculating position, by \SGARefLink{sga3:xxvi:cor:4:4:5}{4.4.5}.  Since \(P\) and \(P_1\) are both
opposite to \(P'_1\), there is a
\[
 g\in\rad^u(P'_1)(S)
\]
such that \(\Int(g)P=P_1\), by \SGARefLink{sga3:xxvi:corollary:4-3-5}{4.3.5}\textup{(i)}.
\SGAQED

For the same reason there is a \(u\in\rad^u(P)(S)\) such that
\(\Int(u)P'=P'_1\).  Hence \(g=\Int(u)u'\) for some
\(u'\in\rad^u(P')(S)\), and
\[
 P_1=\Int(uu'u^{-1})P=\Int(uu')P.
\]
This also gives another proof of \SGARefLink{sga3:xxvi:corollary:5-2}{5.2}.

\SGARefLink{sga3:xxvi:cor:5:3}{Corollaries 5.3} and \SGARefLink{sga3:xxvi:cor:5:4}{5.4} are the essential results of this section.
We first record some consequences of \SGARefLink{sga3:xxvi:cor:5:4}{5.4}.

\begin{sgastatement}{\SGAInlineRefTarget{sga3:xxvi:corollary:5-5}Corollary 5.5}
Let \(S\) be a semilocal scheme and \(G\) an \(S\)-reductive group.
\begin{enumerate}[label=\textup{(\roman*)}]
\item If \(P,Q\) are parabolic subgroups of \(G\) and
\(t(P)\subset t(Q)\), then there is a \(g\in G(S)\) such that
\(\Int(g)P\subset Q\).
\item Let
\[
 P_1\supset P_2\supset\cdots\supset P_n,\qquad
 P'_1\supset P'_2\supset\cdots\supset P'_n
\]
be two chains of parabolic subgroups of \(G\), with
\(t(P_i)=t(P'_i)\) for every \(i\).  There is a \(g\in G(S)\) such
that \(\Int(g)P_i=P'_i\) for every \(i\).
\item Let \(P,Q,P',Q'\) be four parabolic subgroups of \(G\), with
\(t(P)=t(P')\) and \(t(Q)=t(Q')\).  If the pairs \((P,Q)\) and
\((P',Q')\) are both in transverse position, respectively both in
osculating position, then there is a \(g\in G(S)\) such that
\[
 \Int(g)P=P',\qquad\Int(g)Q=Q'.
\]
\item Let \(P,P'\) be parabolic subgroups of the same type, and let
\(L\), respectively \(L'\), be a Levi subgroup of \(P\), respectively
\(P'\).  There is a \(g\in G(S)\) such that
\[
 \Int(g)P=P',\qquad\Int(g)L=L'.
\]
\end{enumerate}
\end{sgastatement}
\SGARefTarget{sga3:xxvi:cor:5:5}

\SGAProof
Assertion \textup{(i)} follows immediately from \SGARefLink{sga3:xxvi:cor:5:4}{5.4}.  Assertion
\textup{(ii)} is proved by induction on \(n\), the case \(n=0\)
being trivial.  We may suppose that \(P_i=P'_i\) for
\(i=1,\ldots,n-1\).  By \SGARefLink{sga3:xxvi:corollary:5-2}{5.2} there is a \(g\in G(S)\) such that
\(\Int(g)P_n=P'_n\).  But \(P_n\) and \(\Int(g)P_n\) are both
contained in \(P_{n-1}=P'_{n-1}\), and hence
\[
 g\in P_{n-1}(S)
\]
by \SGARefLink{sga3:xxvi:lemma:4:4:4}{4.4.4}\textup{(ii)}.  It follows that
\(\Int(g)P_i=P'_i\) for every \(i\).

Assertion \textup{(iv)} follows at once from \SGARefLink{sga3:xxvi:corollary:5-2}{5.2} and \SGARefLink{sga3:xxvi:cor:1:8}{1.8}.  Let us
prove \textup{(iii)} first in the transverse case.  When the types of
\(P\) and \(Q\) are opposite, the assertion follows from
\textup{(iv)} and \SGARefLink{sga3:xxvi:definition:4-3-3}{4.3.3}\textup{(iii)}.  In general, by
\SGARefLink{sga3:xxvi:proposition:4-2-7}{4.2.7}\textup{(iii)} and \SGARefLink{sga3:xxvi:cor:4:4:6}{4.4.6} one can find parabolic subgroups
\[
 P_1\subset P,\quad P'_1\subset P',\quad
 Q_1\subset Q,\quad Q'_1\subset Q'
\]
such that \(P_1,P'_1\) are opposite, \(Q_1,Q'_1\) are opposite, and
\(t(P_1)=t(P'_1)\).  There is consequently a \(g\in G(S)\) such that
\[
 \Int(g)P_1=P'_1,\qquad\Int(g)Q_1=Q'_1.
\]
After conjugating, we may suppose \(P_1=P'_1\) and \(Q_1=Q'_1\).
Then \(P,P'\) are in osculating position and have the same type, so
\(P=P'\) by \SGARefLink{sga3:xxvi:lemma:4:4:4}{4.4.4}\textup{(i)}.  Likewise \(Q=Q'\).

It remains to prove \textup{(iii)} in the osculating case.  By \SGARefLink{sga3:xxvi:corollary:5-2}{5.2}
we may suppose \(P=P'\).  Again by \SGARefLink{sga3:xxvi:corollary:5-2}{5.2} there is a \(g\in G(S)\) such
that \(\Int(g)Q=Q'\).  But then \(g\in P(S)\), by
\SGARefLink{sga3:xxvi:lemma:4:4:4}{4.4.4}\textup{(ii)}, and therefore
\[
 \Int(g)P=P=P'.
\]
\SGAQED
